%!TEX root = ../main.tex

\chapter*{Introduction}\label{ch:introduction}
\markboth{Introduction}{}
\addcontentsline{toc}{chapter}{Introduction}

This work deals with question answering over a collection of documents in a dialogue,
where a user asks several questions in a row and every answer has to come from the
documents. Compared with open-domain \gls{qa}, where a model can use its general
knowledge to answer a single question, this setting adds three difficulties: answers
must be traceable to a specific document, follow-up questions often refer back to
earlier turns (``what about the deadline?''), and some questions cannot be answered from
the documents at all, in which case the system should say so instead of inventing an
answer.

A typical example is an assistant for public services, which has to answer from official
documents and should not make up eligibility conditions or deadlines.

\section*{Motivation}

\Glspl{llm} produce fluent answers to many kinds of questions, but on their own they
can only use what they learned during training. This knowledge is fixed, not specific to
a given set of documents, and leads to hallucinated answers when the model is asked
about details it has never seen.

\Gls{rag}~\cite{lewis2020rag} reduces this problem by retrieving passages from the
documents at question time and giving them to the model together with the question.
The answer can then be based on the actual documents, and the documents can change
without retraining the model. However, the retrieval has to find the right passage
among only a few, the model has to turn several passages into a short answer, and the
system has to notice when none of the passages contains the answer.

In a dialogue there is one more problem: a follow-up question such as ``What about the
deadlines?'' does not contain enough information to retrieve anything by itself. It has
to be rewritten using the earlier turns before retrieval.

\section*{Scope and Research Questions}

The work compares several approaches to this task, starting from a simple baseline and
adding one component at a time. The research questions are:

\begin{enumerate}
    \item How much does \gls{rag} improve over an \gls{llm} used alone, and how much of
          the remaining gap to an oracle that receives the correct passages comes from
          retrieval and how much from generation?

    \item Which retrieval method (sparse, dense or hybrid) works best, and how much does
          rewriting the conversational question help?

    \item Do a cross-encoder reranker and a different generation prompt give significant
          improvements, and where is the bottleneck once retrieval works reasonably well?

    \item Can a knowledge graph built from the same documents by an \gls{llm} be used as
          a retrieval source, on its own or together with passage retrieval?
\end{enumerate}

All experiments use \gls{mdd}~\cite{feng2021multidoc2dial}, a public dataset of
goal-oriented dialogues about US government services, grounded in the pages of four
public websites. Answers are evaluated with token \fone{}, retrieval with Recall@$K$,
and the handling of unanswerable questions with abstention metrics.

\section*{Contributions}

The main results of the work are:

\begin{itemize}
    \item an evaluation setup for \gls{rag} pipelines on \gls{mdd}, with a fixed
          500-example development subset, cached parallel \gls{llm} calls, bootstrap
          confidence intervals and paired significance tests;

    \item a comparison of an \gls{llm}-only baseline with BM25, dense, hybrid and
          rewritten-query retrieval, and an oracle experiment that sets the ceiling for
          generation and separates retrieval errors from generation errors;

    \item the finding that query rewriting is the only change with a clearly
          significant gain in \fone{}, while the reranker improves retrieval without a
          significant gain in the answers, together with an analysis of why;

    \item a knowledge graph of the whole corpus (14{,}731 entities) built with an
          \gls{llm}, with a closed schema derived from a pilot, entity resolution, and
          an evaluation of extraction quality and stability on manually labelled samples;

    \item a comparison of retrieval over the graph with passage retrieval, which shows
          that the graph alone is clearly worse, while facts from the graph added to the
          retrieved passages give a small but significant gain;

    \item a final evaluation of the chosen systems on the full validation split and on
          the held-out test split.
\end{itemize}

\section*{Structure of the Report}

Part~\ref{prt:foundations} gives the background. Chapter~\ref{ch:background} reviews
\glspl{llm}, \gls{rag} and conversational \gls{qa}, and briefly introduces knowledge
graphs. Chapter~\ref{ch:dataset} describes \gls{mdd}, the evaluation metrics and the
dev-500 subset.

Part~\ref{prt:experiments} contains the experiments. Chapter~\ref{ch:llm_baseline}
presents the \gls{llm}-only baseline and Chapter~\ref{ch:rag} the \gls{rag} pipeline
with its retrieval and reranking variants. Chapter~\ref{ch:kg} describes how the
knowledge graph was built and evaluated, and Chapter~\ref{ch:kgret} how it was used for
retrieval. Chapter~\ref{ch:final} takes the remaining design decisions on the full
validation split and reports the final systems on the test split.
Chapter~\ref{ch:conclusions} summarises the findings and the limitations of the work.
